% =====================================================================
%  ch1_intro.tex — 摘要与 §1 引言
% =====================================================================
\begin{abstract}
按通常的表述，亏格 $g\ge 2$ 的曲线的镜像是一个 Landau--Ginzburg 模型；但超势的临界点集才是真正要紧的对象。本文由此出发，对曲线构造一个 A 模型：其对象是三价有理曲线构形，并在节点处带有附加数据。我们为这类三价图定义拉格朗日 Floer 理论与 Fukaya 范畴，并证明它与相应的非阿基米德广义 Tate 曲线的导出范畴等价；随后给出 theta 函数与典范映射的显式对应。
\end{abstract}

\section{引言}\label{s:intro}

黎曼曲面一直是探索同调镜像对称最富饶的例子来源之一：二十年前从椭圆曲线开始 \cite{PZ}，其后又包括一般型代数簇上同调镜像对称最早的若干证据 \cite{SeGenus2,EfimovGenusG,AAEKO}。

要得到任意亏格黎曼曲面的镜像，可以采用多种镜像构造。它们大多关键地依赖于把曲线嵌入某个环绕的环面簇，通常给出三维 Landau--Ginzburg 模型作为镜像，见 \cite{AAK} 的阐述（并见 \cite{HoriVafa,Clarke,ChanLauLeung,GKR}）。不过另有一些构造以带栈结构的节点曲线作为黎曼曲面的镜像 \cite{STZ,GammageShende,LekiliPolishchuk}；在部分情形下，这两类镜像通过 Orlov 广义 Kn\"orrer 周期性的某种形式相联系 \cite{Orlov}。

上述文献所探索的同调镜像对称方向，是把黎曼曲面看作二维辛流形，将其 Fukaya 范畴（A 模型）与镜像 Landau--Ginzburg 模型的奇点导出范畴（B 模型）相比较。本文研究相反的方向：把光滑曲线的凝聚层导出范畴（B 模型）与某个镜像 Landau--Ginzburg 模型的 Fukaya 范畴（A 模型）相比较。这个方向更为困难，部分原因在于：当临界点集非紧时，相应的 Landau--Ginzburg 模型不是恰当的，其 Fukaya 范畴的定义与操作都不容易。在唯一一个镜像 Landau--Ginzburg 模型恰当的例子，即裤子（pair of pants）的情形，可以用微局部层（microlocal sheaf）的语言验证这一等价，见 \cite{Nadler}。关于 $(\C^*)^n$ 中超曲面（本文关心的是 $n=2$ 的情形）这一方向的同调镜像对称，在环面 Landau--Ginzburg 模型的纤维化缠绕 Fukaya 范畴的语言下，已有相当完整的处理，见 \cite{AA}；而亏格 $2$ 曲线嵌入阿贝尔曲面（其 Jacobi 簇）的例子，则用类似方法（除去了紧性带来的问题）在 Cannizzo 的学位论文 \cite{Cannizzo} 中处理。

\cite{AA} 与 \cite{Cannizzo} 所走的路线清楚表明：曲线的 Landau--Ginzburg 镜像，其几何在极大程度上取决于嵌入的选取。事实上，超势的纤维恰好是环绕空间（即曲线所嵌入的空间）的镜像；代数一侧的包含与限制函子 $i_*,i^*$，在镜像对称下对应于 Landau--Ginzburg 模型的 Fukaya 范畴与其正则纤维的 Fukaya 范畴之间的一对伴随函子 $\cup,\cap$。因此，\cite{SeGenus2,GKR,AAK,Cannizzo} 诸文为一族亏格 $2$ 曲线所考虑的种种 Landau--Ginzburg 镜像互不相同，也就不足为奇了：例如 \cite{Cannizzo} 中镜像的奇异纤维不可约，而 \cite{GKR,AAK} 中的奇异纤维有三个不可约分量。然而这些镜像有一个共同特征：在 \cite{SeGenus2} 的情形作 crepant 消解之后，超势的临界点集总由三条有理曲线组成，它们两两相交于两个三重点。类似地，对亏格 $g\ge 2$ 的光滑紧曲线，其镜像超势的临界点集（可能要对全空间作 crepant 消解）由 $3g-3$ 条有理曲线组成，这些曲线相交于 $2g-2$ 个三重点。

就镜像对称的另一方向而言，已有提案主张：可以不再考虑 Landau--Ginzburg 模型的代数几何，而直接考虑该临界点集，并附加使它成为「反常曲线」（perverse curve）的数据 \cite{GKR,Ruddat}。由于奇点导出范畴具有局部性——Orlov 证明了它只依赖临界点集的形式邻域——这一做法一般是合理的。本文的目标是在辛几何（A 模型）一侧做同样的事，从而得到一幅允许显式计算、并且明显不依赖于嵌入选取的曲线同调镜像对称图景。当然这要付出代价：Fukaya--Floer 理论是非局部的，而限制到临界点集会掩盖 Landau--Ginzburg 模型整体辛几何中可能存在的瞬子修正。

我们构造的一般特征由最简单的例子驱动，而这个例子又是其余一切例子的构件：

\begin{figure}[t]
\centering
\begin{tikzpicture}[scale=0.9]
\foreach \a in {0,120,240}{
\begin{scope}[rotate=\a]
  \fill[gray!8] (0.22,-0.38) rectangle (4.75,0.38);
  \draw[gray!50] (0.22,0.38) -- (4.75,0.38);
  \draw[gray!50] (0.22,-0.38) -- (4.75,-0.38);
  \draw[blue,thick] (0,0) -- (4.65,0);
  \draw[cyan,thick,domain=0.3:4.6,samples=170]
    plot (\x,{0.17+0.13*sin(deg(3.1*\x))});
\end{scope}}
\fill (0,0) circle (2.2pt);
\node[above left=1pt and 1pt] at (0,0) {\scriptsize $1$};
\fill (1.9,0) circle (1.2pt);
\node[below=4pt] at (1.9,0) {\scriptsize $t_{1}^{-1}$};
\fill (3.1,0) circle (1.2pt);
\node[below=4pt] at (3.1,0) {\scriptsize $t_{1}^{-2}$};
\node[blue] at (4.1,-0.62) {\scriptsize $L_0$};
\node[cyan] at (4.55,0.78) {\scriptsize $L_0^+$};
\end{tikzpicture}
\caption{裤子镜像 $M=\bigcup_{i=1}^3(\C,0)$ 中的缠绕 Floer 上同调。三条腿各是一个复平面分量，在原点相交成三重点；蓝色直线是拉格朗日图 $L_0$ 的分支，青色曲线示意它的正扰动 $L_0^+$ 沿各端缠绕（图中只标出第一条腿上的生成元 $t_1^{-1},t_1^{-2}$，其余两条腿类推）。}\label{fig:pantsmirror}
\end{figure}

\begin{example}\label{ex:pants}
设 $X$ 为裤子，即 $\PP^1$ 去掉三个点。其镜像 Landau--Ginzburg 模型为 $(\C^3,-xyz)$，临界点集是 $\C^3$ 中三条坐标轴之并。因此我们所取的镜像是构形 $M=\bigcup_{i=1}^3(\C,0)$，即三个复平面 $\C$ 在原点相交于一个三重点。结构层 $\cO_X$ 的镜像是拉格朗日图 $L_0=\bigcup_{i=1}^3\R_{\ge 0}$，由 $M$ 每个分量中的正实轴组成。$L_0$ 在 $M$ 中的缠绕 Floer 上同调有一组加法基：原点处一个生成元，以及 $M$ 的每个端中各一列无穷多个生成元（见图~\ref{fig:pantsmirror}）；它们分别对应于常函数 $1$ 与 $X$ 的三个穿孔点附近坐标 $t_i$ 之逆的逐次幂。然而，一旦考虑 $HW^*(L_0,L_0)$ 上的乘法结构，就立刻看出：$M$ 中拉格朗日 Floer 理论的结构映射必须包含从其中一个分量出发、经原点「传播」（propagate）到另一个分量的全纯圆盘。我们在 §\ref{s:pantsmotivate}--§\ref{s:Amodel} 中进一步解释这一点。
\end{example}

要从裤子过渡到一般情形，先回忆：镜像对称应当在「大复结构极限」附近成立，也就是在非阿基米德框架中成立。Lee 的学位论文 \cite{Lee} 说明了这一普遍预期：曲线的镜像对称与由最大退化诱导的裤子分解相容。确切地说，\cite{AAK} 的构造从 $(\C^*)^2$ 中一列最大退化的复曲线（在热带极限附近）出发，产出环面 Landau--Ginzburg 模型；该镜像由若干标准仿射卡 $(\C^3,-xyz)$ 沿环面坐标变换黏合而成，黏合方式反映了热带极限诱导的曲线组合裤子分解。Lee 构造了该曲线缠绕 Fukaya 范畴的一个版本，可以看作这一裤子分解的 \v Cech 模型，并用它证明了与镜像的奇点导出范畴的等价 \cite{Lee}。

复曲线退化族的语言，在曲线位于镜像对称的辛一侧时是方便的；但在本文的设定中，更富有成效的做法是把曲线 $X$ 定义在非阿基米德域 $K$ 上，即形式变量 $T$ 的、指数为实数的幂级数所成的 Novikov 域——它正是非恰当情形下 Fukaya 范畴的自然定义域。我们考虑的非阿基米德曲线，由光滑化一个最大退化的节点构形 $X^0$ 得到；$X^0$ 是若干带三个标记点的有理曲线之并，各分量按一个三价图两两认同。

\begin{defi}\label{def:combidata}
我们构造所用的\textbf{组合数据}如下。设 $G$ 为有限（无向）图，顶点集为 $V$、边集为 $E$，每个顶点 $v\in V$ 的度数为 $3$，且无自环（即没有从顶点到自身的边）。当 $e\in E$ 与 $v\in V$ 关联时，记作 $e/v$。

对每个 $v\in V$，取 $X^0_v$ 为 $\PP^1_{\Z}$ 的一份拷贝；对每个 $e/v$，固定一个 $\Z$-点 $x_{e/v}\in X^0_v$，使得当 $e\ne e'$ 时 $x_{e/v}$ 与 $x_{e'/v}$ 不相交。

对每个 $e/v$，在 $X^0_v$ 上选取坐标 $t_{e/v}$，使得 $t_{e/v}(x_{e/v})=0$，且 $t_{e/v}$ 在另外两个标记点处取值 $1$ 与 $\infty$。

此外引入形式变量 $\{q_e\}_{e\in E}$，它们将被取为 Novikov 域中赋值 $\val(q_e)=A_e>0$ 的元素。
\end{defi}

在 §\ref{s:Bmodel} 中我们将说明，如何通过光滑化节点曲线
\[
X^0=\Bigl(\bigsqcup_{v\in V}X^0_v\Bigr)\Big/\bigl(x_{e/v}\sim x_{e/v'}\ \ \forall e\in E,\ v\ne v'\bigr)
\]
得到广义 Tate 曲线。用刚性解析几何的话说，这一构造就是把 $X^0$ 的每个节点换成本地坐标下的光滑化 $t_{e/v}t_{e/v'}=q_e$，得到曲线 $X_K$；在 $X_K$ 上，坐标 $t_{e/v}$ 的赋值自然地取值于一个以 $G$ 为模型、边长为 $A_e=\val(q_e)$ 的度量图。

A 一侧是三价构形 $M$，由二维球面组成：其分量与 $E$ 一一对应，节点与 $V$ 一一对应（于是 $M$ 的每个分量都经过两个三重点）。记各分量的辛面积为 $\{A_e\}_{e\in E}$。Fukaya 范畴 $\F(M)$ 在 §\ref{s:Amodel} 中定义。除节点补集中的简单闭曲线之外，该范畴的对象还包括 $M$ 中嵌入的三价图：在每个分量内部，由一条连接该分量两个节点的弧组成。这类对象的 Floer 理论涉及穿过顶点传播的全纯圆盘构形，传播规则由组合数据中的坐标 $t_{e/v}$ 决定（见 §\ref{s:Amodel}）。

我们的主要结果是：

\begin{theo}\label{thm:main}
给定上述组合数据，并令 $q_e=T^{A_e}$，则 Fukaya 范畴 $\F(M)$ 等价于 $\Perf(X_K)$。
\end{theo}
% [Tutor Note: 易错] 注意定理两边的对象都不是「曲线 $X$ 本身」：A 侧是节点构形 $M$（$3g-3$ 条球面粘成），B 侧是广义 Tate 曲线 $X_K$，即 $X$ 的节点退化 $X^0$ 按 $t_{e/v}t_{e/v'}=q_e$ 作非阿基米德光滑化所得。$q_e=T^{A_e}$ 把 B 侧的形式参数与 A 侧的辛面积联系起来，这正是全篇「面积即赋值」的来源。若读者把 $X_K$ 误当成 $X$，就会看不出为何要用 Novikov 域而不用 $\C$。依据：原文 §1 与 §4（本工程 01_通读笔记、03_术语表）。

\begin{remark}\label{rmk:main}
若在 $M$ 上配备 B-场，或对 Fukaya 范畴作\textbf{体形变}（bulk deformation），则本结果可以推广到满足 $\val(q_e)=A_e>0$ 的任意 $q_e\in K$。此外，要求 $G$ 无自环纯粹是为了记号方便，使 $G$ 的半边可以被无歧义地标记；除记号问题之外，本结果及其证明立即推广到带自环的情形。
\end{remark}

\begin{remark}\label{rmk:opencase}
在 A 一侧，我们还可以允许 $M$ 的某些分量为 $S^2\setminus\{pt\}$，即复平面 $\C$，且每个这样的分量上只有一个三重点。这些非紧分量配备无限面积的辛形式，其 Fukaya 范畴既可用无穷远处的缠绕（wrapping）定义，也可用无穷远处的停点（stop）定义。从组合上说，这相当于允许 $G$ 有「外部边」：每个顶点仍连着三条边，但外部边不再连到另一个顶点，也不给外部边配置形式参数 $q_e$。在 B 一侧，若 $x_{e/v}$ 对应一条外部边，则不再把别的分量黏到 $X^0_v$ 的这个标记点上；在缠绕情形，我们从 $X^0$ 与 $X$ 中删去点 $x_{e/v}$；在带停点的情形，则对 $x_{e/v}$ 不作任何处理。例如裤子（例~\ref{ex:pants}）对应于只有一个顶点、三条外部边的情形。本设定下与定理~\ref{thm:main} 相应的结论，可由该定理的证明直接得到。
\end{remark}

\begin{remark}
可以显式验证：仿射椭圆曲线的正则函数环上的乘积结构，与 A 模型上 Floer 乘积的结构常数相符；此例中 A 模型有一个形如 $S^2\setminus\{pt\}$ 的分量（无穷远处取缠绕），以及一个形如 $S^2/(p\sim q)$ 的分量。
\end{remark}

定理~\ref{thm:main} 还有另一推广：考虑接近非最大退化的曲线，即顶点度数可以大于 $3$ 的图。在 B 一侧，这相当于光滑化这样的节点构形：其中每条 $\PP^1$ 可以带三个以上的节点（相应地放宽对构造 $X$ 所用局部坐标 $t_{e/v}$ 的要求）。在 A 一侧，这相当于允许 $M$ 有超过三个分量相接的节点；对象仍然支持在「每个分量内一条连接两节点的弧」所成的图上。定理~\ref{thm:main} 的证明可以适配这一设定，从而在整个刚性解析曲线模空间上建立同调镜像对称。

本文其余部分安排如下。§\ref{s:pantsmotivate} 讨论裤子的情形与 Landau--Ginzburg 模型 $(\C^3,-xyz)$ 的辛几何，以说明我们 A 模型构造的若干关键特征，并指出本构造与其他方法的若干重要差别。§\ref{s:Amodel} 定义我们的 A 模型，即球面三价构形的 Fukaya 范畴。§\ref{s:Bmodel} 说明如何由组合数据构造 B 模型（曲线 $X$），定理~\ref{thm:main} 则在 §\ref{s:equivalence} 中证明；论证用到带 Hamilton 扰动的 Fukaya 范畴 $\F(M)$（类似于 \cite{Lee} 的构造）、同调扰动理论，以及把 $X_K$ 与 $M$ 分解为裤子及其镜像时的限制图。§\ref{ss:theta} 与 §\ref{s:canonicalmap} 展示本框架中 A、B 模型等价的具体程度（与 Landau--Ginzburg 模型的 Fukaya 范畴形成鲜明对照）：我们说明 theta 函数如何由镜像函子的构造产生，并对一般三价图显式确定曲线 $X_K$ 的典范映射及其 A 模型对应物。最后，§\ref{s:higherdim} 试探性地（且高度猜测性地）描述本构造与结果应如何推广到高维。
